\head{Grades K--1: solutions and hints}
\textbf{Plain-words key for the adult.} Stand, sit: circles of 3, 4, 5, 6 leave 3, 1, 3, 5; circles of 7, 8, 9, 10 leave 7, 1, 3, 5; the circle of 16 leaves 1. Stand, stand, sit: circles of 3 to 8 leave 2, 1, 4, 1, 4, 7, and the circle of 9 leaves 1. Use the key to know when to ask for a replay, not to give answers.

\sub{Problem 1, page 1: circles of 3 to 6}
3, 1, 3, 5. Circle of 3: 1 stands, 2 sits, 3 stands, then 1 sits, so 3 is last. Circle of 6: 2, 4, 6 sit; then 1 stands, 3 sits, 5 stands, 1 sits, and 5 is last. A pair that finds 1 last in the circle of 3 has started over at 1. Hint: \say{What word did 3 get? What word comes after stand?}

\sub{Problem 2, pages 2--3: guess, then play 7 to 10}
7, 1, 3, 5. Guesses are spoken; any guess is fine, but ask what it is based on. Children may notice that the last one is never 2, or that the circle of 8 went like the circle of 4. Hint: \say{Which circle that you already played is most like this one?}

\sub{Problem 3, page 4: stars, and the circle of 16}
Stars on the circles of 4 (page 1) and 8 (page 2), and on the circle of 16 once it is played: it leaves 1. A child who says \say{4, 8, then 16} or \say{it keeps doubling} has a real conjecture, and children watching the counters may see why: after the first time round, half are sitting and 1 starts again. Hint: \say{Look at the starred circles. What comes after 4 and 8?} Extension, spoken: which bigger circle would leave 1? (32.)

\sub{Problem 4, page 5: which circles left 3, which left 5}
Number 3: circles of 3, 5 and 9. Number 5: circles of 6 and 10. These lists are complete for the circles on the pages; children find them on their rings. Hint: \say{Where did you draw a ring round a 3?} Extension, spoken: \say{Which circle comes just after a star?} Number 3 comes round again just after each star (5 after 4, 9 after 8), and 5 one circle later; bigger circles 17 and 33 also leave 3, and 18 and 34 leave 5.

\sub{Problem 5, page 5: can 2 be last?}
No. Number 2 always gets the first \say{sit}, in every circle that has a 2. \say{It never happened} is evidence, not yet a reason; the follow-up \say{Could it happen in a circle we haven't played?} invites the reason. Hint: \say{In your circles, when did 2 sit down?} Extension: can 4 be last? (No: 4 gets the second \say{sit}; every even number sits the first time round.)

\sub{Problem 6, pages 6--7: stand, stand, sit, circles of 3 to 8}
2, 1, 4, 1, 4, 7. Circle of 3: 3 sits; then 1 and 2 stand and 1 sits, so 2 is last. Circle of 4: 3 sits; 4 and 1 stand, 2 sits; 4 and 1 stand, 4 sits; 1 is last. The three-word chant across the wrap is harder than before, so the finger-and-flip roles matter more. Hint: \say{Where is your finger now, and who gets the next word?}

\sub{Problem 7, page 8: where to stand in the circle of 9}
At 1. The sitting order is 3, 6, 9, 4, 8, 5, 2, 7. Choosing where to stand is the same as finding who is last; children may guess, then play to find out. Hint: \say{How could you find out without guessing?}

\sub{Problem 8, page 8: can 3 be last with stand, stand, sit?}
No. Number 3 always gets the first \say{sit}. Hint: \say{In your stand, stand, sit circles, when did 3 sit down?} Extension: can 6 be last? (No: 6 gets the second \say{sit}.)

\head{Grades 2--3: solutions and hints}
\sub{Problem 1, page 1: circles of 1 to 16}
1, 1, 3, 1, 3, 5, 7, 1, 3, 5, 7, 9, 11, 13, 15, 1. The circle of 1 has nobody to sit. The pattern: a 1 at 1, 2, 4, 8 and 16, with the odd numbers climbing by 2 in between. Hint: \say{Where are the 1s? What happens just after each one?}

\sub{Problem 2, page 1: predict 20, 25, 32}
9, 19, 1. Continuing the climb after 16: 17 leaves 3, \dots, 20 leaves 9, 25 leaves 19, 31 leaves 31, and 32 leaves 1. Predictions go in the table before playing. Each check takes 19 to 31 sits on the 32-spot mat; if a play disagrees with a prediction, replay before deciding which is wrong. Hint: \say{What did your table do after 8? What might it do after 16?}

\sub{Problem 3, page 1: numbers that can never be last}
Exactly the even numbers. The first time round everyone is standing, so the words go stand, sit, stand, sit in order, and every even number sits before the chant reaches 1 again. Every odd number does win somewhere: each climb after a power of two runs through all the odd numbers below the next power of two (17 is last in the circle of 24). For a child who answers from the table alone: \say{Could an even number win in a circle bigger than 16?}

\sub{Problem 4, page 2: why 16 and 32 leave 1}
In the circle of 16, the first lap seats 2, 4, \dots, 16. Number 16 got the lap's last word, a \say{sit}, so 1 gets \say{stand} again: it is now a circle of 8 starting at 1. The next laps seat 3, 7, 11, 15, then 5, 13, then 9. Each lap halves the circle and ends on a \say{sit}, so 1 is never told to sit. After the first lap of 32, the 16 odd numbers remain with 1 next, which is the circle of 16 again. On the small ring children can cross out one lap at a time. Hints: \say{After the first time round, how many are standing, and who says the next word?} \say{Have you played a circle like that?}

\sub{Problem 5, page 2: number 9, and the same number}
Number 9: circles of 12, 20 and 36. Same number: 1, 3, 7, 15, 31 (accept with or without 1). The lists are complete because 9 appears once in each climb that reaches it, four circles after the power of two 8, 16 or 32; the climb after 4 stops at 7, and the one after 64 starts beyond 40. The last one equals the circle size at the top of each climb, just before a power of two. Hint: \say{Where does 9 come in each climb?} Discretionary follow-up: \say{How do you know there are no more up to 40?}

\sub{Problem 6, page 3: circles of 50 and 100}
37 and 73. The nearest power of two below 50 is 32, which leaves 1; each extra person adds 2, and 50 is 18 more, so $1+2\times18=37$. Another reason: once 2, 4, \dots, 36 have sat (18 people), 32 remain and 37 is next, a fresh circle of 32 starting at 37. For 100: $100=64+36$, so $1+72=73$. Hints: \say{Which circle near 50 are you sure about?} \say{What happens to the answer when one more person joins?} Why the climb goes up by 2 is the step rule in the overview.

\sub{Problem 7, page 3: every third, circles of 1 to 12}
1, 2, 2, 1, 4, 1, 4, 7, 1, 4, 7, 10. In the circle of 2 the words go 1 stand, 2 stand, 1 sit, so 2 is last. Hint for small circles: \say{What if you keep going round, even with only two people?}

\sub{Problem 8, page 4: predict 13 to 20, check 14 and 20}
13, 2, 5, 8, 11, 14, 17, 20; the checks give 2 and 20 (20 needs the 32-spot mat). The pattern: add 3 to the previous answer, and if that is more than the circle size, take the circle size away ($10+3=13$; $13+3=16$ and $16-14=2$). The 1--12 table drops to 1 at 4, 6 and 9, so a child may predict \say{back to 1}; the circle of 14 lands on 2, which is why it is a check. Hint: \say{How much does it go up each time? What happens when it would go past the circle size?}

\sub{Problem 9, page 4: where to stand among 30}
At 29. Continuing from 20, circles of 21 to 30 leave 2, 5, 8, 11, 14, 17, 20, 23, 26, 29. Playing 30 on the 32-spot mat is a long check, not the main route. Hint: \say{Can your table keep going past 20?} Extension: Grades 4--5 Problem 9 (41 people; stand at 31).

\head{Grades 4--5: solutions and hints}
\sub{Problems 1 and 2, page 1}
As Grades 2--3 Problems 1 and 2: the table 1, 1, 3, 1, 3, 5, 7, 1, 3, 5, 7, 9, 11, 13, 15, 1, and 20, 25, 32 leave 9, 19, 1.

\sub{Problem 3, page 1: every circle up to 100 that leaves 1}
1, 2, 4, 8, 16, 32, 64 (accept without 1). The first lap explains both directions. In an even circle the first lap seats the even numbers and 1 starts again in a circle half the size; in an odd circle bigger than 1, the last person of the first lap stands and then 1 sits. So 1 is last exactly when halving never meets an odd number above 1, which means a power of two. The small ring is for crossing out the laps of 16. Hint: \say{After the first lap, how many are left, and who goes next?} Discretionary follow-up: \say{How do you know no other circle up to 100 leaves 1?}

\sub{Problem 4, page 2: circles of 100 and 1000}
73 and 977. $100=64+36$: once 36 people have sat (2, 4, \dots, 72), 64 remain and 73 is next, so by Problem 3's reasoning 73 is last. $1000=512+488$: after 2, 4, \dots, 976 sit, 977 starts a circle of 512. Hint: \say{Which power of two is just below 100? How many must sit to get down to it, and who is next?}

\sub{Problem 5, page 2: from 26 to 52 and 53}
41 and 43. In the circle of 52, after the first lap $1,3,\dots,51$ remain with 1 next. That is the circle of 26 with person $j$ renamed $2j-1$, so the 21st is $2\times21-1=41$. In the circle of 53, the first lap seats 2 to 52, then 53 stands and 1 sits; $3,5,\dots,53$ remain with 3 next, person $j$ renamed $2j+1$, so the answer is $2\times21+1=43$. Hint: \say{After the first lap of 52, who is still standing, and who goes next? What circle does that look like?} Extension: the general doubling rule, and the chain 25, 50, 100 (19, 37, 73) as a check on Problem 4.

\sub{Problem 6, page 3: number 23, number 22}
Number 23: circles of 27, 43 and 75. Number 22: none, because every even number sits in the first lap. The 23 list is complete because $23=2\times11+1$: the circle must be a power of two plus 11, with the power bigger than 11, so $16+11$, $32+11$ and $64+11$ ($8+11$ fails, since the climb after 8 stops at 15, and $128+11$ is over 100). In children's terms, 23 comes 11 circles after a power of two, if that climb is long enough. Hints: \say{How many circles after 16 does the climb reach 23?} For 22: \say{When does 22 sit down?}

\sub{Problem 7, page 3: every third, circles of 1 to 12}
As Grades 2--3 Problem 7: 1, 2, 2, 1, 4, 1, 4, 7, 1, 4, 7, 10.

\sub{Problem 8, page 4: predict 13 to 24, check 14 and 21}
13, 2, 5, 8, 11, 14, 17, 20, 2, 5, 8, 11; both checks give 2 (21 needs the 32-spot mat). Both land just past a wrap, so they test \say{add 3, and take away the circle size} against \say{drop to 1}. For a child who asks why the step works: once 3 sits, 4 starts a circle of $n-1$ (the proof is the step rule's, in the overview). Hint: \say{How does each answer compare with the one before?}

\sub{Problem 9, page 4: where to stand among 41}
At 31. Continuing the step rule from 24: circles of 25 to 30 leave 14, 17, 20, 23, 26, 29; the circle of 31 leaves 1; then 4, 7, 10, \dots, and the circle of 41 leaves 31. Number 16 is the last to sit, so 16 and 31 are the last two standing. Hint: \say{Do you need to play 41, or can the table take you there?}

\sub{Problem 10, page 5: the binary rule}
Move the leading 1 to the end, dropping any zeros that end up in front: $13=1101\to1011=11$, $25=11001\to10011=19$, $16=10000\to00001=1$. Why: the leading 1 stands for the largest power of two $2^m$ in $n$, and the remaining digits are $L$. Moving the 1 to the end shifts $L$'s digits up one place, which doubles it, and adds 1, giving $2L+1$, the rule behind Problem 4's answers. The printed example rules out \say{reverse the digits} ($1100\to0011=3$). \say{Swap the second and last digits} fits 12 but fails at once on 8 and 16. Hint: \say{Try your rule on 8 and on 16. What should it give?} If binary is new: \say{Which of 32, 16, 8, 4, 2, 1 add up to 25?}

\head{Sources, evidence and what to record}
Ronald L. Graham, Donald E. Knuth and Oren Patashnik, \emph{Concrete Mathematics}, 2nd ed. (Addison-Wesley, 1994), \S1.3, \say{The Josephus Problem}, which the outline cites for the recurrences and the binary rotation. Not consulted for this guide (no copy is in the project's references); page numbers unchecked. The proofs above are written independently.

Year 1, Handouts 6.7 and 7.8 (\texttt{lowell-math-circle-year-1}): consulted. They ask how to order a deck so that a deal gives $1,2,\dots,n$, but do not state the deal. If it moves every other card to the bottom, it is this elimination run backward; confirm with the organizer before calling this week a return.

Evidence: \texttt{check\_guide.py} (in this guide's source folder) simulates the game literally, confirms every statement in the overview for every second person up to $n=2048$ and every third up to 500, and recomputes every answer, list and count printed here, including the launch and its replay on counters, the flip counts behind the timings, and the spot spacing. These are digital checks; nothing here has been rehearsed with counters or children.

Library connections, from the outline: Week 4 hops a fixed number of dots round a ring, Week 89 grows patterns layer by layer, and Weeks 8 and 30 use binary and ternary. New here: an elimination process on a circle with a doubling recurrence and a binary rule.

Record after a session: preparation time; whether the counters fit the spots and the mat halves stayed together; whether pairs kept the chant going past 1 without reminders; which entries were wrong and how they were caught; which problems each table reached; and what children wanted to continue.
